% ============================================================
%  A Ground-Based Telemetry Analysis Pipeline for CubeSat
%  Anomaly Detection with Synthetic Fault Injection
%  IEEE Conference Paper — IEEEtran two-column format
% ============================================================
\documentclass[conference]{IEEEtran}

\usepackage[T1]{fontenc}
\usepackage[utf8]{inputenc}
\usepackage{amsmath}
\usepackage{graphicx}
\usepackage{booktabs}
\usepackage{siunitx}
\usepackage{caption}
\usepackage{subcaption}
\usepackage{hyperref}
\usepackage{xcolor}
\usepackage{tabularx}
\usepackage{multirow}

\hypersetup{
  colorlinks=true,
  linkcolor=blue,
  citecolor=blue,
  urlcolor=blue
}

% ── Title ────────────────────────────────────────────────────
\title{A Ground-Based Telemetry Analysis Pipeline for CubeSat
       Anomaly Detection with Synthetic Fault Injection}

\author{
  \IEEEauthorblockN{Logan Luna}
  \IEEEauthorblockA{
    \textit{KISPE Space Systems / University Collaboration}\\
    loganluna@example.com
  }
}

\begin{document}

\maketitle

% ── Abstract ─────────────────────────────────────────────────
\begin{abstract}
Reliable anomaly detection in spacecraft telemetry is essential for
mission assurance, yet labelled fault datasets from real hardware
tests remain scarce.  This paper presents a modular ground-based
telemetry analysis pipeline developed for the KISPE Satellite
Learning Laboratory~(SLL), a CubeSat-class hardware-in-the-loop
platform.  The pipeline automates parsing of SCOTTI archive files,
cleans and decodes binary telemetry fields, applies four complementary
anomaly detectors — linear thermal trend analysis, z-score flagging,
cross-channel correlation breakdown, and stuck-sensor detection —
and generates up to twenty diagnostic plots per experiment session.
Applied to nine completed subsystem test campaigns (totalling 3{,}340
CDH and 3{,}392 ADCS telemetry packets), the pipeline consistently
identifies a universal on-board data handling~(OBDH) thermal ramp of
\SIrange{0.8}{2.0}{\celsius\per\minute} across all sessions, 21
thermally-stuck sensor channels, and experiment-specific
voltage-rail anomalies.  To support future machine-learning
classifier development, a synthetic fault injection framework
generates 16 labelled anomaly types — spanning thermal, power, ADCS,
reaction wheel, and timing fault classes — producing 71 labelled CSV
artefacts per experiment from real baseline telemetry.  The full
pipeline, injector, and dataset are released openly alongside this
paper.
\end{abstract}

\begin{IEEEkeywords}
CubeSat, anomaly detection, telemetry analysis, fault injection,
satellite health monitoring, ADCS, synthetic dataset generation
\end{IEEEkeywords}

% ============================================================
\section{Introduction}
% ============================================================

Spacecraft anomaly detection has traditionally relied on threshold
monitors and expert-defined rules embedded in ground software.
While effective for known, well-characterised faults, such approaches
struggle to generalise to novel failure modes and require extensive
manual configuration for each new mission.  The emergence of
machine-learning~(ML) anomaly detection — in particular unsupervised
and semi-supervised methods — offers a path toward more adaptive,
data-driven fault identification~\cite{hundman2018detecting,
kerner2020comparison}.  A fundamental bottleneck, however, is the
availability of labelled fault datasets from real satellite hardware;
most operational anomaly datasets are proprietary, sparsely
labelled, or gathered only in simulation.

The KISPE Satellite Learning Laboratory~(SLL) addresses this gap from
the hardware perspective.  It is a laboratory-scale, CubeSat-derived
platform that provides a fully integrated Command and Data Handling
(CDH) and Attitude Determination and Control System~(ADCS) in a
representative flight-like configuration, with real-time telemetry
streamed via the SCOTTI ground interface.  In this work we exploit the
SLL to conduct nine distinct subsystem test campaigns and to develop a
complete software pipeline for:

\begin{enumerate}
  \item \textbf{Telemetry parsing} — ingesting both Analysis-format and
        Debug-format SCOTTI archives, with automatic fallback;
  \item \textbf{Data cleaning} — decoding signed/unsigned hex fields,
        deduplicating packets, deriving elapsed time, and flagging
        stuck sensor channels;
  \item \textbf{Anomaly detection} — applying four complementary
        statistical detectors that operate on CDH and ADCS streams;
  \item \textbf{Visualisation} — generating up to 20 diagnostic plots
        per session for exploratory analysis;
  \item \textbf{Synthetic fault injection} — producing a labelled,
        ML-ready dataset by injecting 16 fault types into the cleaned
        real-world baseline telemetry.
\end{enumerate}

The contributions of this paper are:
\begin{itemize}
  \item A reproducible, open-source ground-test telemetry pipeline
        for CubeSat-class hardware;
  \item A characterisation of naturally occurring anomalies observed
        across nine SLL test campaigns;
  \item A systematic fault injection taxonomy and implementation
        covering thermal, power, ADCS, wheel, and timing fault classes;
  \item A labelled dataset of 71 CSV artefacts per experiment,
        derived from real hardware telemetry, suitable for training and
        evaluating anomaly detection models.
\end{itemize}

The remainder of the paper is structured as follows.
Section~\ref{sec:background} reviews related work.
Section~\ref{sec:platform} describes the SLL platform and telemetry
format.
Section~\ref{sec:pipeline} details the pipeline architecture.
Section~\ref{sec:experiments} describes the test campaigns.
Section~\ref{sec:results} presents detection results.
Section~\ref{sec:injection} describes the synthetic fault injection
framework.
Section~\ref{sec:discussion} discusses findings and limitations, and
Section~\ref{sec:conclusion} concludes.

% ============================================================
\section{Background and Related Work}
\label{sec:background}
% ============================================================

\subsection{Spacecraft Anomaly Detection}

Early spacecraft health monitoring relied on hard-coded limit checks
applied to individual telemetry channels~\cite{iverson2008general}.
Subsequent work introduced correlation-based monitors, which flag
deviations in pairs of channels that are normally tightly
coupled~\cite{srivastava2005discovering}.  Statistical process control
methods, including Shewhart charts and exponentially weighted moving
averages, have been applied to satellite voltage and temperature
streams~\cite{susto2018multivariate}.

\subsection{Machine Learning Approaches}

Long Short-Term Memory~(LSTM) networks have shown particular promise
for multivariate spacecraft telemetry anomaly detection, with the
SMAP and MSL Mars rover datasets becoming de-facto
benchmarks~\cite{hundman2018detecting}.  Telemanom and related
prediction-error methods flag time steps where a trained sequence
model cannot reconstruct the observed signal.  More recently,
transformer-based architectures~\cite{xu2021anomaly} have been
applied to satellite telemetry, though they demand larger labelled
corpora than are typically available from single missions.

\subsection{Synthetic Data and Fault Injection}

The lack of labelled spacecraft fault data has motivated several
synthetic generation strategies.  \cite{kerner2020comparison}
benchmarked statistical and deep methods on simulated Mars rover
data.  Fault injection at the software level — inserting error
patterns directly into clean sensor streams — has been used in
avionics and industrial process control to create balanced training
sets~\cite{natschlaeger2010fault}.  To the authors' knowledge, no
prior work has applied a systematic, multi-class fault injection
framework to real CubeSat hardware telemetry at the subsystem
test-campaign level.

% ============================================================
\section{Platform and Telemetry Format}
\label{sec:platform}
% ============================================================

\subsection{KISPE Satellite Learning Laboratory}

The KISPE SLL is a CubeSat-derived educational and research platform
comprising an On-Board Computer~(OBC), an Electrical Power System~(EPS)
with battery and solar simulation, a Communications module~(COMMS), and
a full ADCS board hosting a 3-axis accelerometer, gyroscope,
magnetometer, sun sensors, a reaction wheel, and magnetorquers.
The platform generates two telemetry streams at approximately
\SI{1}{\hertz}--\SI{2}{\hertz}:

\begin{description}
  \item[CDH packets] contain power rail voltages and currents
        (SYS\_V, PLAT\_5V\_V, OBC\_PWR\_V, COMMS\_PWR\_V, \ldots),
        nine board temperatures (TEMP\_SOLAR through TEMP\_PAYLOAD),
        20 expansion thermistor channels, and housekeeping counters.
  \item[ADCS Sensor packets] carry 3-axis inertial measurements
        (ACCEL, GYRO, MAG), eight sun sensor channels, sun angle,
        wheel speed, and magnetorquer commands.
\end{description}

\subsection{SCOTTI Archive Format}

The SCOTTI ground interface writes telemetry to two file types per
session:

\textbf{Analysis files} contain full decoded fields in human-readable
key-value format (\texttt{FIELD\_NAME = value}), one packet per
timestamped block, with sub-millisecond timestamps.

\textbf{Debug files} record raw hex payloads together with
space-delimited ``Plot'' lines that carry all 87 CDH and 30 ADCS
decoded values in fixed positional order.  Debug files use
\texttt{MM/DD/YYYY} date formatting and \SI{1}{\second} timestamp
precision, compared to the \texttt{DD/MM/YYYY} sub-millisecond format
of Analysis files.  Experiments where the Analysis file was empty
(Magnetorquer\_1, AttitudeDeterminationMagnetometer\_1,
AttitudeDetermination\_Control, CameraTest\_1) were ingested via the
Debug fallback.

% ============================================================
\section{Pipeline Architecture}
\label{sec:pipeline}
% ============================================================

The pipeline is implemented in Python~3 using pandas, NumPy, SciPy,
and Matplotlib, and is structured into five modules:
\texttt{parser}, \texttt{cleaner}, \texttt{detector},
\texttt{plotter}, and \texttt{injector}.
Fig.~\ref{fig:pipeline} illustrates the data flow.

\subsection{Parsing (\texttt{parser.py})}

The parser applies a regular-expression state machine to scan
sequential lines of an archive file.  For Analysis files, a
timestamp-header regex identifies each packet boundary, and a
tab-indented field regex extracts name-value pairs.  For Debug files,
a single line regex matches Plot lines and maps the space-delimited
tokens to 87 CDH or 30 ADCS field names using fixed positional
schemas.  Hex strings that cannot be converted to float are preserved
as strings for downstream handling.  The parser returns two
\texttt{pd.DataFrame} objects (CDH and ADCS), sorted by timestamp.

\subsection{Cleaning (\texttt{cleaner.py})}

Cleaning applies several transformations:
\begin{enumerate}
  \item \textbf{Hex decoding}: unsigned-integer fields (current
        readings) are converted via \texttt{int(str, 16)}; signed
        16-bit fields (BATT\_CHR\_I, WHEEL\_SPEED, MAGNETORQUER\_*) use
        two's-complement decoding.  Mixed string/float columns produced
        by duplicate Debug sessions are handled by checking the input
        type before attempting hex parsing.
  \item \textbf{Deduplication}: rows with identical timestamps
        (\texttt{keep=`last'}) are removed.  This eliminates the
        duplicate sessions present in several Debug files.
  \item \textbf{Elapsed time}: a scalar \texttt{elapsed\_s} column is
        derived from the packet timestamp minus the session start.
  \item \textbf{Power estimation}: estimated power in milliwatts is
        computed for four supply rails (SYS, PLAT\_5V, OBC\_PWR,
        COMMS\_PWR) from voltage $\times$ decoded current.
  \item \textbf{Stuck thermistor detection}: any thermistor channel
        with zero variance across all packets is labelled stuck.
\end{enumerate}

\subsection{Detection (\texttt{detector.py})}

Four detectors are applied to each cleaned session:

\textbf{D1 — Linear thermal trend.}  A least-squares line is fitted
to TEMP\_OBDH and TEMP\_COMMS vs.\ elapsed time using
\texttt{scipy.stats.linregress}.  The slope (°C/min), $R^2$, and
total temperature rise $\Delta T$ are reported.

\textbf{D2 — Z-score flagging.}  For each of seven CDH channels
(TEMP\_OBDH, TEMP\_COMMS, TEMP\_SOLAR, SYS\_V, PLAT\_5V\_V,
OBC\_PWR\_V, COMMS\_PWR\_V), samples with $|z| > 2.5\sigma$ are
flagged, where $z$ is computed over all clean samples in the session.

\textbf{D3 — Correlation breakdown.}  A rolling Pearson correlation
($w = 30$ samples) is computed for four board-temperature pairs that
are physically expected to track closely:
(TEMP\_EPS, TEMP\_BATTERY),
(TEMP\_EPS, TEMP\_BACKPLANE),
(TEMP\_COMMS, TEMP\_PAYLOAD), and
(TEMP\_ADCS, TEMP\_WHEEL).
Windows where $|r| < 0.5$ are flagged as correlation breakdowns.

\textbf{D4 — Stuck-channel detection.}  Any ADCS sensor channel with
zero variance but a non-zero stuck value (excluding the all-zero
baseline state) is identified as a frozen sensor.

\subsection{Plotting (\texttt{plotter.py})}

Up to 20 plots are generated per session.  Plots 01–15 target CDH
data (board temperatures, OBDH ramp, power rails, thermistor
anomaly, correlation matrices, voltage overlays, etc.).  Plots 16–20
are conditional on live ADCS data and cover IMU 3-axis time-series,
vector magnitudes, reaction wheel speed, sun sensor readings, and
magnetorquer commands.

% ============================================================
\section{Experimental Campaigns}
\label{sec:experiments}
% ============================================================

Nine test campaigns were conducted on the KISPE SLL between
02 March and 22 March 2026, each lasting between approximately
\SI{3}{\minute} and \SI{10}{\minute}.
Table~\ref{tab:experiments} summarises the dataset.

\begin{table}[t]
\caption{SLL Test Campaign Overview}
\label{tab:experiments}
\centering
\footnotesize
\begin{tabular}{@{}llrrc@{}}
\toprule
\textbf{Experiment} & \textbf{Source} & \textbf{CDH} & \textbf{ADCS} & \textbf{ADCS Live} \\
\midrule
Baseline\_1                       & Analysis & 333  & 335  & No  \\
Accelerometer\_1                  & Analysis & 360  & 363  & Yes \\
Gyro\_1                           & Analysis & 267  & 264  & Yes \\
ReactionWheel\_1                  & Analysis & 357  & 359  & Yes \\
SunSensorTest\_1                  & Analysis & 361  & 361  & Yes \\
Magnetorquer\_1                   & Debug    & 535  & 534  & Yes \\
AttitudeDet.\ Magnetometer\_1     & Debug    & 285  & 285  & Yes \\
AttitudeDetermination\_Control    & Debug    & 273  & 273  & Yes \\
CameraTest\_1                     & Debug    & 545  & 546  & Yes \\
\midrule
\textbf{Total}                    &          & \textbf{3{,}316} & \textbf{3{,}280} & \\
\bottomrule
\end{tabular}
\end{table}

The Baseline\_1 session was conducted with ADCS powered off as a
control experiment.  Subsequent campaigns powered the full ADCS stack
and exercised individual actuators and sensors in isolation.  The
CameraTest\_1 and AttitudeDetermination sessions were the longest
(approximately \SI{9}{\minute}) and were captured only in Debug
format due to a logging configuration change on 22 March 2026.

% ============================================================
\section{Results}
\label{sec:results}
% ============================================================

\subsection{Universal OBDH Thermal Ramp}

The most prominent system-wide finding is a consistent linear rise in
the OBDH board temperature across every experiment.
Table~\ref{tab:obdh} reports the fitted slope, goodness-of-fit $R^2$,
and total temperature delta for all nine sessions.

\begin{table}[t]
\caption{OBDH Thermal Ramp Characteristics per Session}
\label{tab:obdh}
\centering
\footnotesize
\begin{tabular}{@{}lrrr@{}}
\toprule
\textbf{Experiment} & \textbf{Slope} & $\boldsymbol{R^2}$ & $\boldsymbol{\Delta T}$ \\
                    & (\si{\celsius\per\minute}) &  & (\si{\celsius}) \\
\midrule
Baseline\_1                    & 1.52 & 0.975 & 4.48 \\
Accelerometer\_1               & 1.92 & 0.945 & 6.09 \\
Gyro\_1                        & 2.04 & 0.948 & 4.85 \\
ReactionWheel\_1               & 1.56 & 0.943 & 5.24 \\
SunSensorTest\_1               & 1.26 & 0.944 & 4.24 \\
Magnetorquer\_1                & 0.80 & 0.888 & 8.98 \\
AttitudeDet.\ Magnetometer\_1  & 1.25 & 0.921 & 6.97 \\
AttitudeDetermination\_Control & 1.18 & 0.933 & 6.06 \\
CameraTest\_1                  & 1.06 & 0.956 & 10.66 \\
\midrule
\textbf{Mean $\pm$ SD} & $1.40 \pm 0.39$ & $0.939 \pm 0.024$ & $6.51 \pm 2.06$ \\
\bottomrule
\end{tabular}
\end{table}

All sessions exhibit $R^2 > 0.88$, confirming highly linear thermal
behaviour.  The Gyro\_1 session shows the steepest rate
(\SI{2.04}{\celsius\per\minute}), while Magnetorquer\_1 has the
lowest slope but highest total $\Delta T$ (\SI{8.98}{\celsius}) owing
to its longer duration.  CameraTest\_1 achieves a $\Delta T$ of
\SI{10.66}{\celsius} — the largest observed — consistent with the
payload being powered throughout the \SI{9}{\minute} session.

This ramp is not an injected anomaly; it reflects steady-state power
dissipation in the OBC in the absence of active thermal control, and
serves as a reliable discriminant between operational and
anomalous thermal states.

\subsection{Stuck Thermistor Array}

All experiments register 21 stuck thermistor channels: 16 rod
thermistors frozen at \SI{77.18}{\celsius} and four bath/expansion
thermistors frozen at \SI{74.31}{\celsius}, with the heater
thermistor reading zero.  These values are characteristic of an
unconnected or open-circuit ADC input at the conversion resolution
of this firmware version, and represent a systemic sensor interface
fault rather than a transient event.  CameraTest\_1 is the sole
exception, with only one stuck channel, suggesting a firmware or
cabling update between the March 12 and March 22 campaigns.

\subsection{Voltage Rail Z-Score Anomalies}

Table~\ref{tab:zscore} summarises the z-score anomaly counts for
the four primary voltage rails.  SYS\_V and OBC\_PWR\_V show
elevated flag counts in most sessions, reflecting sub-cycle
transients on the battery-backed supply rail — a known
characteristic of the SLL's charge-controller interaction at low
solar input.  PLAT\_5V\_V anomalies are highest in the Gyro\_1 and
ReactionWheel\_1 sessions, consistent with increased current demand
from the ADCS peripherals.

\begin{table}[t]
\caption{Z-Score Anomaly Counts ($|z|>2.5$) for Key Channels}
\label{tab:zscore}
\centering
\footnotesize
\begin{tabular}{@{}lrrrr@{}}
\toprule
\textbf{Experiment} & \textbf{SYS\_V} & \textbf{PLAT\_5V} & \textbf{OBC\_PWR} & \textbf{COMMS} \\
\midrule
Baseline\_1                    &  3 &  3 &  9 &  0 \\
Accelerometer\_1               & 17 & 19 & 17 &  2 \\
Gyro\_1                        & 18 & 24 &  3 & 19 \\
ReactionWheel\_1               &  7 & 20 & 14 & 23 \\
SunSensorTest\_1               &  6 &  3 &  2 &  4 \\
Magnetorquer\_1                & 16 &  9 & 17 &  2 \\
AttitudeDet.\ Magnetometer\_1  & 21 & 13 &  5 &  9 \\
AttitudeDetermination\_Control & 12 &  0 &  3 &  0 \\
CameraTest\_1                  &  3 &  0 &  0 &  0 \\
\bottomrule
\end{tabular}
\end{table}

\subsection{ADCS Sensor State}

The Baseline\_1 session shows all nine IMU channels
(ACCEL\_X/Y/Z, GYRO\_X/Y/Z, MAG\_X/Y/Z) identically zero,
confirming that the ADCS board was unpowered.  All subsequent
sessions show live ADCS data.  Two frozen-sensor findings are
notable:

\begin{itemize}
  \item \textbf{SUN\_ANGLE frozen at 0x013C~(316)} — detected in the
        Accelerometer\_1 and Magnetorquer\_1 sessions.  The sun angle
        derivation algorithm was likely not re-initialised between the
        two sessions, causing the output register to hold the last
        computed value.
  \item \textbf{WHEEL\_SPEED frozen near zero} — observed in Gyro\_1.
        The reaction wheel was not commanded during this session and
        the speed register did not update, indicating a firmware
        reporting oversight rather than a mechanical fault.
\end{itemize}

\subsection{Temperature Correlation Breakdown}

The rolling-correlation detector flags the
(TEMP\_EPS, TEMP\_BACKPLANE), (TEMP\_COMMS, TEMP\_PAYLOAD), and
(TEMP\_ADCS, TEMP\_WHEEL) pairs in all nine sessions.  This pattern
is consistent across experiments, suggesting that these pairs are
not, in fact, tightly correlated on the SLL time-scale of a few
minutes — the boards heat at slightly different rates due to their
varying power dissipation.  This is a useful baseline characterisation:
a true thermal decoupling anomaly (e.g.\ a heater failure) would need
to be distinguished from this normal intra-session decorrelation.

\subsection{Packet Timing}

Sessions parsed from Analysis files exhibit a mean CDH inter-packet
interval of $\approx \SI{0.50}{\second}$ with a standard deviation of
$\approx \SI{0.14}{\second}$, reflecting the 2 Hz nominal telemetry
cadence.  Debug-format sessions show exactly $\SI{1.00}{\second}$
intervals with zero standard deviation, an artefact of the Debug
format's integer-second timestamp resolution.  No pathological packet
gaps exceeding \SI{2}{\second} were observed.

% ============================================================
\section{Synthetic Fault Injection}
\label{sec:injection}
% ============================================================

\subsection{Motivation and Design Principles}

Real spacecraft fault datasets are intrinsically imbalanced: nominal
operation dominates, and most fault types have very few or zero
observed instances.  To address this for ML classifier training, we
develop a fault injection framework operating directly on the cleaned
real-world telemetry.  This preserves the realistic noise
characteristics, sensor correlations, and thermal drift of the actual
hardware while adding controlled, labelled perturbations.

Each injection function is a pure function of its input
\texttt{pd.DataFrame}, returning a modified copy alongside a label
record containing the anomaly type, start/end row indices, severity
(\textit{low}/\textit{medium}/\textit{high}), and affected channel
names.  All random decisions use a seeded \texttt{numpy.Generator},
making every injection fully reproducible.

\subsection{Injection Taxonomy}

Table~\ref{tab:injections} lists the 16 implemented injection types
across five fault classes.

\begin{table}[t]
\caption{Synthetic Fault Injection Taxonomy}
\label{tab:injections}
\centering
\footnotesize
\begin{tabular}{@{}llp{4.2cm}@{}}
\toprule
\textbf{ID} & \textbf{Name} & \textbf{Description} \\
\midrule
\multicolumn{3}{@{}l}{\textit{Thermal (CDH)}} \\
T1 & OBDH ramp accel.      & Accelerate OBDH slope by $\times$4 from a mid-session split \\
T2 & Board step change      & Step $\delta$T on a board temp with exponential decay ($\tau=30$ samples) \\
T3 & Board divergence       & Grow a linear bias between two normally-correlated boards \\
\midrule
\multicolumn{3}{@{}l}{\textit{Power (CDH)}} \\
P1 & SYS\_V sag             & Gaussian-shaped battery sag; co-dips BATT\_CHR\_V \\
P2 & OBC voltage spikes     & 8 random-sign spikes in a 20-sample OBC\_PWR\_V window \\
P3 & Rail stuck high        & Force ADCS\_5V\_V to 3.3 V from a random mid-session row \\
\midrule
\multicolumn{3}{@{}l}{\textit{ADCS (ADCS stream)}} \\
A1 & Accel axis dropout     & Zero ACCEL axes for a contiguous window \\
A2 & Gyro saturation        & Peg a gyro axis at its observed maximum \\
A3 & Magnetometer flip      & Invert all MAG axes for a short burst \\
A4 & IMU all-zero dropout   & Zero all 9 IMU channels for 25\% of session \\
A5 & Sun sensor saturation  & Set all SUN\_SENSOR channels to 1023 \\
\midrule
\multicolumn{3}{@{}l}{\textit{Wheel (ADCS stream)}} \\
W1 & Wheel overspeed        & Linearly ramp WHEEL\_SPEED to $\pm$1800 counts \\
W2 & Wheel sudden stop      & Zero WHEEL\_SPEED one sample after peak \\
\midrule
\multicolumn{3}{@{}l}{\textit{Timing (CDH)}} \\
C1 & Packet gap             & Remove rows totalling 10 s from mid-session \\
\midrule
\multicolumn{3}{@{}l}{\textit{Compound}} \\
CP1 & Thermal + power fault & T1 + P1 applied simultaneously \\
CP2 & ADCS + power fault    & A4 on ADCS + P3 on CDH simultaneously \\
\bottomrule
\end{tabular}
\end{table}

\subsection{Dataset Statistics}

For each of the nine processed experiments, two variants of each
injection type are generated using different random seeds, yielding
$16 \times 2 = 32$ injected data files and 32 corresponding label
CSVs per experiment.  Additionally, two clean evaluation copies
(with \texttt{anomaly=False} labels) and one \texttt{injection\_summary.csv}
are written, bringing the total to 71 artefacts per experiment and
$71 \times 9 = 639$ artefacts across the full campaign dataset.
Each label record encodes the anomaly type, affected channels,
severity level, and exact row span, providing rich supervision signal
for both binary and multi-class fault classifiers.

% ============================================================
\section{Discussion}
\label{sec:discussion}
% ============================================================

\subsection{Key Findings}

The universal OBDH thermal ramp (mean
\SI{1.40}{\celsius\per\minute}, $R^2 > 0.88$ in all sessions) is the
dominant feature of SLL telemetry and can serve as a reliable
normality baseline: any session that does not exhibit this trend would
be immediately anomalous.  The distinct rates across experiments
(\SIrange{0.80}{2.04}{\celsius\per\minute}) suggest that subsystem
power draw significantly modulates the ramp slope, with the gyro test
showing the steepest rise due to the sustained ADCS peripheral load.

The 21 stuck thermistors in 8 of 9 sessions indicate an unresolved
hardware or firmware issue with the expansion thermistor interface on
the March 12 hardware configuration.  The fact that CameraTest\_1
(22 March) shows only one stuck channel suggests the fault was
partially corrected between campaigns.  These channels are flagged
deterministically and excluded from statistical analyses.

The voltage rail anomalies (Table~\ref{tab:zscore}) appear correlated
with ADCS peripheral load rather than representing true power faults.
This contextual dependency is an important consideration for ML
classifier design: a model trained on Baseline\_1 (low z-score counts)
may over-flag normal ADCS-on behaviour in subsequent sessions.  The
synthetic injection dataset mitigates this by generating fault
examples from the relevant baseline sessions.

\subsection{Limitations}

Several limitations constrain the scope of the current study:

\begin{enumerate}
  \item \textbf{Session duration}: all sessions are 3–10 minutes,
        limiting the ability to characterise slow-drift anomalies
        that manifest over longer timescales.
  \item \textbf{Missing ThermalTest\_1}: both files for this session
        are empty (\SI{1}{line}), representing a logging capture
        failure.  Re-running the session is recommended.
  \item \textbf{Debug timestamp resolution}: second-level timestamps
        in Debug files preclude packet-level timing analysis and
        result in $\sim$50\% row loss after deduplication for some
        sessions.
  \item \textbf{Injector realism}: synthetic injections are applied
        as mathematical transformations and do not account for
        secondary physical effects (e.g.\ a power sag would also
        affect downstream temperature trajectories).  Higher-fidelity
        injection would require a coupled physics model.
\end{enumerate}

\subsection{Future Work}

Immediate next steps include: (i) training and evaluating LSTM and
Isolation Forest classifiers on the labelled injection dataset;
(ii) extending session durations to 30–60 minutes to capture thermal
equilibrium behaviour; (iii) implementing a higher-resolution Debug
file parser that disambiguates sub-second packets using the firmware
SECOND counter field; and (iv) adding a detection module for
magnetorquer / reaction wheel coupling anomalies using the
wheel--gyro rolling correlation already computed by the pipeline.

% ============================================================
\section{Conclusion}
\label{sec:conclusion}
% ============================================================

This paper presented a complete telemetry analysis pipeline for the
KISPE Satellite Learning Laboratory, applied to nine subsystem test
campaigns.  The pipeline ingests both Analysis and Debug SCOTTI archive
formats, applies four complementary anomaly detectors, and generates
rich diagnostic visualisations — all fully automated.  Key empirical
findings include a universal OBDH thermal ramp of
\SI{1.40}{\celsius\per\minute} on average, a systemic 21-channel
thermistor interface fault, experiment-dependent voltage transients
correlated with ADCS peripheral loading, and two frozen ADCS sensor
registers.  To support data-driven fault detection research, a
16-class synthetic fault injector produces 639 labelled telemetry
artefacts across all campaigns from real hardware baseline data.
The pipeline and dataset are made openly available.

% ── Bibliography ─────────────────────────────────────────────
\bibliographystyle{IEEEtran}
\bibliography{references}

\end{document}
